
\section{AT-02: All terms for two equivalent d electrons}
\source{2026 Atomic Assignment 3, Question 2; Week 8 angular-momentum lectures.}
\subsection*{Step 1: count allowed antisymmetric states}
A \(d\) electron has \(l=2\), with five \(m_l\) values
and two spin projections: 10 spin-orbitals. Two equivalent
electrons cannot occupy the same spin-orbital,
so there are
\[
\boxed{N_{\rm microstates}=\binom{10}{2}=45}
\]
independent antisymmetric states. This count must equal
the total degeneracy of the final terms.

\subsection*{Step 2: couple spatial and spin angular momenta}
The orbital total can be
\(L=0,1,2,3,4\); two spins \(s=1/2\) give
\(S=0\) (antisymmetric singlet) or \(S=1\)
(symmetric triplet). For two identical \(l=2\)
electrons, exchanging their positions in the
coupled orbital state contributes \((-1)^{2l-L}=(-1)^L\).
Fermionic antisymmetry requires:
\[
\begin{array}{c|c|l}
\text{Orbital exchange}&\text{Spin exchange}&\text{Allowed }L,S\\\hline
L\text{ even (symmetric)}&S=0\text{ (antisymmetric)}&
{}^1S,\ {}^1D,\ {}^1G\\
L\text{ odd (antisymmetric)}&S=1\text{ (symmetric)}&
{}^3P,\ {}^3F
\end{array}
\]

\subsection*{Step 3: attach the total electronic \(J\)}
For each term, \(J=|L-S|,\ldots,L+S\):
\[
\boxed{\begin{gathered}
{}^1S_0,\quad {}^1D_2,\quad{}^1G_4,\\
{}^3P_{0,1,2},\quad{}^3F_{2,3,4}.
\end{gathered}}
\]
Check all degeneracies:
\[
\underbrace{1}_{^1S}+\underbrace5_{^1D}
+\underbrace9_{^1G}+\underbrace{(1+3+5)}_{^3P}
+\underbrace{(5+7+9)}_{^3F}
=\boxed{45}.
\]

\subsection*{Step 4: apply Hund's rules with their caveats}
First maximise total spin \(S\), favouring triplets;
next maximise \(L\), favouring the \({}^3F\) term;
since \(d^2\) is less than half full, lowest \(J\)
is favoured within that term:
\[
\boxed{\text{Predicted ground term: }{}^3F_2}.
\]
These rules do not uniquely order \emph{all}
excited terms. Consult the specified species in
the NIST Atomic Spectra Database for the measured order.
\checkmethod{Only listing \(L=0,\ldots,4\) and \(S=0,1\)
overcounts forbidden states. Exchange antisymmetry and
the 45-microstate check are essential.}
